\chapter{Cơ sở lý thuyết và nghiên cứu liên quan}
\label{chap:background}

\section{Agent dựa trên mô hình ngôn ngữ lớn}
Mô hình ngôn ngữ lớn có thể được dùng như một ``agent'' khi nhận đầu vào là quan sát về môi trường, suy luận bằng ngôn
ngữ tự nhiên và trả về một hành động. Kỹ thuật chuỗi suy luận (chain-of-thought)~\cite{wei2022cot} yêu cầu mô hình
viết ra các bước trung gian trước khi kết luận, giúp cải thiện khả năng suy luận nhiều bước và cho phép kiểm tra lại
lập luận. Reflexion~\cite{shinn2023reflexion} bổ sung khả năng tự rút kinh nghiệm: agent đọc lại kết quả các hành động
trước đó và viết nhận xét bằng ngôn ngữ để điều chỉnh hành vi ở các bước sau. Cơ chế này được CryptoTrade sử dụng dưới
dạng agent phản tư (reflection) xem lại kết quả giao dịch của 7 ngày trước.

Một hạn chế cố hữu của agent LLM là hiện tượng ``ảo giác'': mô hình có thể đưa ra dữ kiện không có trong đầu vào. Trong
bài toán tài chính, đây là rủi ro nghiêm trọng, vì một con số bịa đặt có thể dẫn tới quyết định sai. Do đó hệ thống
trong đề tài yêu cầu mọi lập luận trong giai đoạn tranh biện phải trích dẫn dòng bằng chứng cụ thể và được kiểm tra
trước khi được chấp nhận (Mục~\ref{sec:debate}).

\section{Hệ đa tác tử và tranh biện}
\paragraph{Tranh biện đa tác tử.}
Du và cộng sự~\cite{du2024debate} đề xuất cho nhiều phiên bản LLM cùng trả lời một câu hỏi, sau đó mỗi phiên bản đọc
câu trả lời của các phiên bản khác và cập nhật câu trả lời của mình qua nhiều vòng. Phương pháp này cải thiện độ chính
xác trên các bài toán suy luận và giảm dữ kiện sai. Liang và cộng sự~\cite{liang2024divergent} chỉ ra hiện tượng ``thoái hóa
suy nghĩ'' (Degeneration-of-Thought) khi một LLM tự phản tư, và dùng tranh biện giữa các agent có lập trường đối lập cùng một trọng tài để
khuyến khích tư duy phân kỳ. ChatEval~\cite{chan2024chateval} áp dụng tranh biện để nhiều agent cùng đánh giá chất
lượng văn bản, cho thấy việc gán vai trò khác nhau cho các agent là yếu tố quan trọng.

\paragraph{Tranh biện khi các nguồn mâu thuẫn.}
MADAM-RAG~\cite{wang2025madamrag} xét bài toán hỏi đáp dựa trên truy xuất khi các tài liệu truy xuất được mâu thuẫn
nhau, chứa thông tin sai hoặc nhiễu. Mỗi tài liệu được giao cho một agent riêng, các agent tranh biện qua nhiều vòng
và một bộ tổng hợp tóm tắt cuộc tranh biện để đưa ra câu trả lời cuối. Ý tưởng mỗi nguồn một agent, tranh biện khi
mâu thuẫn và tổng hợp lịch sử tranh biện qua các vòng là nguồn cảm hứng trực tiếp cho kiến trúc của đề tài.

\paragraph{Tranh biện trong quyết định tài chính.}
Huang và cộng sự~\cite{huang2026debate} cho các agent chuyên môn đề xuất, phản biện và sửa quyết định phân bổ danh mục
trong 210 lần mô phỏng trên dữ liệu thị trường lịch sử. Chất lượng lập luận của tranh biện (tính hợp lệ, bằng chứng,
cân nhắc phương án khác) không có quan hệ với chỉ số Sharpe hay lợi suất. Nhóm tác giả chỉ ra lỗi ``hội tụ xu nịnh''
(sycophantic convergence): qua các vòng phản biện, agent bỏ quan điểm độc lập và hội tụ về cùng một lựa chọn. Một biện
pháp giữ lại bất đồng giữa các agent (dựa trên phân kỳ Jensen--Shannon) làm tăng Sharpe, cho thấy tranh biện có giá trị
khi nó \emph{giữ} được các tín hiệu độc lập chứ không phải khi nó làm lập luận đẹp hơn.

\section{Agent LLM cho giao dịch tài chính}
FinMem~\cite{yu2023finmem} thiết kế agent giao dịch cổ phiếu có bộ nhớ phân tầng theo tầm thời gian. FinCon~\cite{yu2024fincon}
tổ chức nhiều agent theo mô hình quản lý -- nhân viên và dùng phản hồi bằng ngôn ngữ để cập nhật niềm tin đầu tư.
TradingAgents~\cite{xiao2024tradingagents} mô phỏng một công ty giao dịch với các agent phân tích cơ bản, tâm lý, tin
tức, kỹ thuật và hai nhóm nghiên cứu lạc quan -- bi quan tranh luận trước khi agent giao dịch quyết định.

Với tiền mã hóa, CryptoTrade~\cite{li2024cryptotrade} là agent LLM kết hợp dữ liệu on-chain và tin tức cho giao dịch
không cần huấn luyện (zero-shot), đánh giá trên BTC, ETH và SOL trong ba giai đoạn thị trường năm 2023. Luo và
cộng sự~\cite{luo2025cryptomas} xây dựng hệ đa tác tử gồm các nhóm agent chuyên trách dữ liệu thị trường và tin tức cho
quản lý danh mục 30 đồng tiền mã hóa, có so sánh các cấu trúc giao tiếp phân cấp, cộng tác và tranh biện. Gần đây,
Market Regime Council~\cite{mrc2026} đề xuất gán trọng số động cho ba agent chuyên môn dựa trên giá trị Shapley tính từ
lịch sử hiệu quả, kèm hệ số nhân theo trạng thái thị trường.

Gần nhất với đề tài là CGX~\cite{bajgora2026cgx} (Consensus-Gated Execution): hai agent Bull và Bear tranh biện ba vòng,
một Meta-Evaluator chấm mức đồng thuận và chỉ cho phép vào lệnh khi đồng thuận đủ mạnh, mặc định là đứng ngoài thị
trường. Đánh giá trên tiền mã hóa giai đoạn 2022--2025, lợi ích chính được báo cáo là giảm mức sụt giảm vốn tối đa và độ
biến động. Khác với đề tài, các agent của CGX được chia theo \emph{phe} (lạc quan -- bi quan) chứ không theo \emph{nguồn
dữ liệu}, mức đồng thuận do mô hình chấm chứ không đo bằng công thức, lập luận không được kiểm tra nguồn, và baseline là
các chiến lược đơn giản chứ không phải một agent LLM đã công bố được cài đặt lại.

\section{Baseline CryptoTrade}
\label{sec:cryptotrade}
CryptoTrade gồm bốn thành phần LLM chạy mỗi ngày:
\begin{itemize}
  \item \textbf{Agent phân tích thị trường và on-chain}: đọc giá mở cửa gần đây, ba tín hiệu kỹ thuật (giao cắt đường
  trung bình ngắn -- dài, MACD, dải Bollinger) và các thống kê giao dịch on-chain.
  \item \textbf{Agent phân tích tin tức}: đọc tối đa 5 bài báo về đồng tiền trong ngày trước đó.
  \item \textbf{Agent phản tư}: xem lại các quyết định và lợi nhuận của 7 ngày trước.
  \item \textbf{Agent giao dịch}: đọc ba bản phân tích trên và trả về một hành động trong $[-1, 1]$; giá trị dương là
  tỷ lệ tiền mặt dùng để mua, giá trị âm là tỷ lệ coin đem bán.
\end{itemize}
Điểm then chốt là \emph{mọi mâu thuẫn giữa các nguồn được agent giao dịch xử lý ngầm trong một lời nhắc}. Bài báo
có thí nghiệm loại bỏ từng nguồn, nhưng chỉ báo cáo cho ETH trong giai đoạn tăng giá; phiên bản chỉ dùng dữ liệu thị trường cho BTC trong đề tài
do đề tài tự chạy.

\section{Các công cụ định lượng}
\paragraph{Phân kỳ Kullback--Leibler.} Với hai phân phối rời rạc $P$ và $Q$ trên cùng tập kết quả, phân kỳ
KL~\cite{kullback1951}
\begin{equation}
  D_{KL}(P \parallel Q) = \sum_{j} P(j)\log\frac{P(j)}{Q(j)}
\end{equation}
đo lượng thông tin mất đi khi dùng $Q$ để xấp xỉ $P$. $D_{KL} \geq 0$, bằng 0 khi và chỉ khi $P = Q$, và không đối
xứng. Trong đề tài, phân kỳ KL được dùng để đo khoảng cách giữa niềm tin của hai agent.

\paragraph{Phương sai quyết định.} Phương sai của các giá trị niềm tin có dấu của các agent đo độ phân tán trực tiếp của
quyết định, bổ sung cho phân kỳ KL vốn nhạy với các phân phối gần tất định.

\paragraph{Chỉ báo kỹ thuật.} Đường trung bình động hàm mũ (EMA) làm mượt chuỗi giá với trọng số giảm dần theo thời
gian. Chỉ số sức mạnh tương đối RSI trong khoảng $[0,100]$ đo tỷ lệ giữa mức tăng và mức giảm trung bình trong 14 phiên;
RSI gần 70 trở lên thường được xem là quá mua, gần 30 trở xuống là quá bán. MACD là hiệu giữa EMA 12 và EMA 26 so với
đường tín hiệu EMA 9 của chính nó. Dải Bollinger gồm đường trung bình 20 phiên và hai dải cách nó hai độ lệch chuẩn.

\paragraph{Chỉ số Sharpe.} Chỉ số Sharpe~\cite{sharpe1966} đo lợi suất trên một đơn vị rủi ro. Theo đúng cách tính
của CryptoTrade, với $r_t$ là lợi suất ngày và lãi suất phi rủi ro bằng 0:
\begin{equation}
  \text{Sharpe} = \frac{\overline{r}}{\sigma(r)}.
\end{equation}

\section{Khoảng trống nghiên cứu}
Các hệ thống giao dịch đa tác tử hiện có hoặc gộp các nguồn trong một lời nhắc (CryptoTrade), hoặc dùng tranh luận như
một bước sinh lập luận mà không đo mức mâu thuẫn một cách định lượng (TradingAgents). CGX dùng mức đồng thuận để chặn
lệnh, nhưng đồng thuận giữa hai phe chứ không giữa các nguồn dữ liệu; Huang và cộng sự phân tích tranh biện trên thị
trường mô phỏng. Các nghiên cứu về tranh biện đa tác tử còn lại chủ yếu đánh giá trên bài toán hỏi đáp có đáp án đúng
xác định. Chưa có đánh giá có kiểm soát nào so sánh việc xử lý mâu thuẫn tường minh \emph{giữa các nguồn dữ liệu} với
cách gộp ngầm của một baseline đã công bố, trên cùng dữ liệu thật và cùng mô hình ngôn ngữ, trong bài toán giao dịch
tiền mã hóa. Đề tài hướng tới lấp khoảng trống này.
